\chapter{}

\section{\lecon{22.2} 阿拉伯语书法}

\begin{quote}
    默德那……其书旁行，兼篆、草、楷三体，西洋诸国皆用之。
    
    \hfill ——《明史·西域四》
\end{quote}

\begin{itemize}
    \item \ac{اَلْمَدِينَةُ}{麦迪那（默德那）}
    \item \ac{اَلْكُوفَةُ}{库法} 
    \item \ac{اَلْخَطُّ اَلْكُوفِيُّ}{库法体}
    \item \ac{خَطُّ النَّسْخِ}{誊抄体（纳斯赫体）}
    \item \ac{خَطَّ الثُّلُثِ}{三一体（苏鲁斯体）}
    \item \ac{اَلْخَطُّ الْفَارِسِيُّ}{波斯体}
    \item \ac{اَلْخَطُّ الدِّيوَانِي}{公文体（迪瓦尼体）}
    \item \ac{خَطُّ الرُّقْعَةِ}{行书体（卢格阿体）}
    \item \ac{بِسْمِ اللهِ الرَّحْمٰنِ الرَّحِيمِ}{奉至仁至慈的真主之名（太斯米）}
    \item \ac{لَا إِلَهَ إِلَّا اللهُ مُحَمَّدٌ رَسُولُ اللهِ}{（清真言）} 
    \item \ac{فَكِّرْ كَثِيرًا وَتَكَلَّمْ قَلِيلًا، تَكَلَّمْ قَلِيلًا وَافْعَلْ كَثِيرًا}{多思少言，少言多行}
    \item \ac{اَلْمَغْرِبُ}{马格里布}
    \item \ac{اَلْخَطُّ الْمَغْرِبِيُّ}{马格里布体}
    \item \ac{اَلْأَنْدَلُسُ}{安达卢斯（今西班牙附近）}
    \item \ac{اَلْخَطُّ الصِّينِيُّ}{中国体}
    \item \ac{اَلطُّغْرَاءُ}{花押}
\end{itemize}

\begin{attention}
    马格里布体中，\arm{فـ} 写作 \arm{ڢـ}，\arm{قـ} 写作 \arm{ڧـ}。
\end{attention}

\begin{note}
    太斯米中，课上给出的写法是 \arm{الرَّحْمَن} 但特意强调这里的 \arm{ـمَـ} 读长音，参照下文中的太斯米独立编码的写法，此处直接记为 \arm{الرَّحْمٰنِ} 。
\end{note}

% \newfontfamily\ruq[Script=Arabic]{ArefRuqaa-Regular}
% \begin{Arabic}
%     \Large
%     {\ruq فكّر كثيرًا وتكلّم قليلًا، تكلَّم قليلًا وافعل كثيرًا}
%     % 字体测试 تــــــكلَّم قليلًا وافعل كثيرًا https://mehbarah.com/index.html
% \end{Arabic}

\newpage
\begin{note}
    以下不是课程内容。

    \subsection*{计算机字库的支持}
    
    课堂上是用图片形式展示各类书法的，但现在的字库已经对各类字体做了诸多适配，不展示一下就太可惜了。以下针对课堂提到的书法字体，各选了一种能正确显示的计算机字体展示。选用标准是：字库名中必须明确指出属于哪种字体。

    \newfontfamily\kuf[Script=Arabic]{DiwanKufi}
    \newfontfamily\nas[Script=Arabic]{AdobeNaskh-Medium}
    \newfontfamily\thu[Script=Arabic]{AThuluth}
    \newfontfamily\far[Script=Arabic]{Farisi}
    \newfontfamily\diw[Script=Arabic]{DiwaniLetter}
    \newfontfamily\ruq[Script=Arabic]{ArefRuqaa-Regular}
    \newfontfamily\magh[Script=Arabic]{AALMAGHRIBI}

    \begin{center}
        \renewcommand{\arraystretch}{1.8}
        \begin{tabular}{c|r|c}
            字体 & \multicolumn{1}{c|}{示例文字} & 字库名\\
            \hline 
            库法 & \arm{\rm \large \kuf اطلب العلم ولو في الصين} & Diwan Kufi\\
            纳斯赫 & \arm{\rm \Large \nas اطلب العلم ولو في الصين} & Adobe Naskh\\
            苏鲁斯 & \arm{\rm \Large \thu اطلب العلم ولو في الصين} & A Thuluth\\
            波斯 & \arm{\rm \Large \far اطلب العلم ولو في الصين} & Farisi\\
            迪瓦尼 & \arm{\rm \Large \diw اطلب العلم ولو في الصين} & Diwani Letter\\
            卢格阿 & \arm{\rm \Large \ruq اطلب العلم ولو في الصين} & Aref Ruqaa\\
            马格里布 & \arm{\rm \Large \magh اطلب العلم ولو في الصين} & AALMAGHRIBI\\
        \end{tabular}
        \renewcommand{\arraystretch}{1}
    \end{center}
    
    \subsection*{太斯米的特殊支持}
    
    一些计算机字体支持特殊的太斯米 Unicode 合字字符，特展示出来。注意，每个案例都只有一个 Unicode 字符（U+FDFD），由 \LaTeX 直接编译而成。

    \begin{center}
        % \newfontfamily\tsmx{Alkalami-Regular}
        % {\rm \large \tsmx ﷽}\\
        % （字库：Alkalami）

        % ~\\
        \newfontfamily\tsmxi{Ruwudu-Regular}
        {\rm \LARGE \tsmxi ﷽}\nopagebreak\\
        （字库：Ruwudu）

        ~\\
        \newfontfamily\tsmvii{IBMPlexSansArabic-Regular}
        {\rm \Large \tsmvii ﷽}\nopagebreak\\
        （字库：IBM Plex Sans Arabic）

        % ~\\
        % \newfontfamily\tsmv{Firjar}
        % {\rm \Large \tsmv ﷽}\nopagebreak\\
        % （字库：Firjar）

        ~\\
        \newfontfamily\tsmix{HONORSansArabicUI-R}
        {\rm \Large \tsmix ﷽}\nopagebreak\\
        （字库：HONOR Sans Arabic UI）

        ~\\
        \newfontfamily\tsmi{NotoNaskhArabic-Regular}
        {\rm \Huge \tsmi ﷽}\nopagebreak\\
        （字库：Noto Naskh Arabic）

        ~\\
        \newfontfamily\tsmii{Damascus}
        {\rm \Huge \tsmii ﷽}\nopagebreak\\
        （字库：Damascus）

        ~\\
        \newfontfamily\tsmviii{NotoSansArabic-Regular}
        {\rm \Huge \tsmviii ﷽}\nopagebreak\\
        （字库：Noto Sans Arabic）

        ~\\
        \newfontfamily\tsmiv{ScheherazadeNew-Regular}
        {\rm \Huge \tsmiv ﷽}\nopagebreak\\
        （字库：Scheherazade New）

        ~\\
        \newfontfamily\tsmxii{Almarai-Regular}
        {\rm \Huge \tsmxii ﷽}\nopagebreak\\
        （字库：Almarai）

        ~\\
        \newfontfamily\tsmvi{Gulzar}
        {\rm \huge \tsmvi ﷽}\nopagebreak\\
        （字库：Gulzar）

        % ~\\
        % \newfontfamily\tsmiii{NotoNastaliqUrdu}
        % {\rm \Large \tsmiii ﷽}\nopagebreak\\
        % （字库：Noto Nastaliq Urdu）

        ~\\
    \end{center}
    
\end{note}
